\documentclass[10pt,letterpaper]{article}
\usepackage{cloudsuite}
\renewcommand{\DocShort}{CH. 07}
\begin{document}
\coverpage{Chapter 7: Understanding Cloud Security and Cybersecurity}{Cornell study notes and applied review | FOUNDATIONS}{Concepts \textbullet{} decisions \textbullet{} verification artifacts}
\sourceStatement
\tableofcontents
\newpage
\section{Learning contract}
\begin{takeaway}
Build a threat model grounded in trust boundaries and responsibility allocation.
\end{takeaway}
\textbf{Chapter scope:} Security and threat vocabulary, actors, attacks and risk management.\\
\sourceRef{7}{159}
\subsection{Mapped source sections}
\begin{itemize}
\item 7.1 Basic Security Terminology
\item 7.2 Basic Threat Terminology
\item 7.3 Threat Agents
\item 7.4 Common Threats
\item 7.5 Case Study Example
\item 7.6 Additional Considerations
\item 7.7 Case Study Example
\end{itemize}
\section{Cornell cue and notes matrix}
\small\setlength{\LTleft}{0pt}\setlength{\LTright}{0pt}
\begin{longtable}{>{\RaggedRight\arraybackslash}p{0.29\textwidth}>{\RaggedRight\arraybackslash}p{0.63\textwidth}}
\toprule\textbf{Cue / essential question}&\textbf{Explanatory notes and reasoning}\\\midrule\endfirsthead\toprule\textbf{Cue}&\textbf{Notes}\\\midrule\endhead
\textbf{How are security qualities distinguished?} & Confidentiality concerns disclosure, integrity concerns unauthorized change, availability concerns access when needed, and authenticity concerns trust in identity or origin.\\[3pt]\textcolor{petrol}{\emph{Practice:} Associate one control and evidence item with each quality.}\\[8pt]\midrule
\textbf{What is the difference between threat and risk?} & A threat describes a potential adverse event or actor capability; risk evaluates expected impact and likelihood under actual exposure and controls.\\[3pt]\textcolor{petrol}{\emph{Practice:} Rank two threats using asset criticality and current exposure.}\\[8pt]\midrule
\textbf{Why do overlapping trust boundaries matter?} & Cloud stakeholders and technical layers can have distinct administrative authority. Failing to resolve a boundary permits unintended privilege or data access.\\[3pt]\textcolor{petrol}{\emph{Practice:} Place trust crossing points on a multitenant service diagram.}\\[8pt]\midrule
\textbf{Which threat families must be modeled?} & Consider eavesdropping, malicious intermediaries, denial of service, authorization faults, virtualization or container attacks, malicious code and privilege escalation.\\[3pt]\textcolor{petrol}{\emph{Practice:} Map three attack paths to preventive and detective controls.}\\[8pt]\midrule
\textbf{How can insider and social attacks affect cloud control planes?} & Credentials, trusted access and human workflow decisions can bypass purely technical perimeter measures.\\[3pt]\textcolor{petrol}{\emph{Practice:} Define a privileged access review and phishing-resistant login control.}\\[8pt]\midrule
\textbf{What makes a security program actionable?} & Policy, contract, monitoring, risk owner and remediation evidence must align. A technology control that lacks operations can silently become ineffective.\\[3pt]\textcolor{petrol}{\emph{Practice:} Build a threat-to-control-to-test traceability row.}\\[8pt]\midrule
\end{longtable}\normalsize
\section{Synthesis and architecture connections}
Translate the chapter vocabulary into one explicit service requirement, one boundary, one owner and one measurable test. Later mechanism and architecture chapters provide possible realizations; the requirements remain the acceptance baseline.
\begin{workbox}[title={Applied deliverable}]
Build a threat model grounded in trust boundaries and responsibility allocation.\\[4pt]Create a one-page architecture or operations artifact with: (1) assumptions, (2) a decision and alternatives, (3) control ownership, (4) observable acceptance criteria, and (5) a failed-case test.
\end{workbox}
\section{Active recall}
\begin{enumerate}
\item How are security qualities distinguished?
\item What is the difference between threat and risk?
\item Why do overlapping trust boundaries matter?
\item Which threat families must be modeled?
\end{enumerate}
\subsection{Answer-check rubric}
\begin{itemize}
\item A correct response distinguishes the named concepts rather than relying on vendor names.
\item The explanation identifies an operational consequence, boundary or tradeoff.
\item At least one measurement, test or documentary artifact could demonstrate the claim.
\end{itemize}
\section{Blank Cornell capture sheet}
\textbf{Topic / source subsection:}\enspace\rule{0.65\linewidth}{0.4pt}\\[6pt]
\begin{longtable}{p{0.28\textwidth}p{0.64\textwidth}}
\toprule Cues / questions & Notes / evidence \\ \midrule
\rule{0pt}{42pt} & \\ \midrule
\rule{0pt}{42pt} & \\ \midrule
\rule{0pt}{42pt} & \\ \midrule
\rule{0pt}{42pt} & \\ \midrule
\rule{0pt}{42pt} & \\ \midrule
\end{longtable}
\textbf{Summary in three sentences:}\\[16pt]\hrulefill\\[16pt]\hrulefill\\[16pt]\hrulefill

\refs
\end{document}
